\documentclass[11pt,letterpaper]{article}
\usepackage[utf8]{inputenc}
\usepackage[T1]{fontenc}
\usepackage[spanish]{babel}
\usepackage{amsmath,amssymb}
\usepackage[margin=1in]{geometry}
\usepackage{enumitem}
\usepackage{booktabs}
\usepackage{pgfplots}
\pgfplotsset{compat=1.17}
\setlength{\parindent}{0pt}
\newcommand{\kt}{\tilde{k}}

\begin{document}

\begin{center}
  {\Large\bfseries Comentes de convergencia -- PEP 1}\\[2pt]
  {\large\bfseries Macroeconomía II}\\[6pt]
  \rule{\textwidth}{0.4pt}\\[2pt]
  Casos reales de África (datos del Banco Mundial) · Repaso del 6-11 y clase 4
\end{center}

\begin{itemize}
  \item En los \textbf{comentes}, indique si la afirmación es Verdadera, Falsa o Incierta, y justifique con ecuaciones,
  gráficos y conceptos (convergencia absoluta versus condicional, dinámica de transición, efecto nivel versus
  efecto crecimiento). Una respuesta sin justificación no es una respuesta.
  \item Use los datos de la tabla y el gráfico. La pauta está al final: no la mire hasta terminar.
\end{itemize}

\section*{Lo que dijo el profesor en el repaso}
Fuera del estado estacionario la economía crece (o decrece), y \emph{mientras más lejos esté, más rápido}:
\[
\frac{\dot{k}}{k} = s\,\frac{y}{k} - (n+\delta)
\qquad\text{(con tecnología: }\ \frac{\dot{\kt}}{\kt} = s\,\kt^{\,\alpha-1} - (n+g+\delta),\ \ \kt = k/A\text{)}.
\]
\textbf{Convergencia absoluta:} todos tienen la misma función de producción y el mismo estado estacionario, así que
los pobres crecen más rápido y la desigualdad desaparece sola.
\textbf{Convergencia condicional:} cada país converge a \emph{su propio} estado estacionario,
\[
y_i^*(t) = A_i(t)\,h_i\left(\frac{s_i}{n_i+g+\delta}\right)^{\frac{\alpha}{1-\alpha}},
\]
y crece más rápido cuanto más lejos está \emph{de ese}.

\section*{Datos}
PIB por persona en dólares constantes de 2015 (Banco Mundial, \texttt{NY.GDP.PCAP.KD}). $s$ se aproxima con la
formación bruta de capital (\% del PIB) y $n$ con el crecimiento de la población.

\begin{center}
\footnotesize
\setlength{\tabcolsep}{4pt}
\begin{tabular}{lrrlrl}
\toprule
País & $y$ 1960 & $y$ 2024 & Crec. anual de $y$ & $n$ 2000-23 & Dato clave \\
\midrule
Botsuana & 394 & 6.952 & 9,0\% (66-90); 1,3\% (00-24) & 1,7\% & $s \approx 35\%$ (66-90) \\
RD del Congo & 1.303 & 554 & $-1{,}3\%$ (60-24) & 3,2\% & $s \approx 12\%$ (90-10) \\
Ruanda & 333 & 1.043 & 8,3\% (94-04) & 2,2\% & PIB: $-50\%$ (1994), $+35\%$ (1995) \\
Níger & 742 & 589 & $-0{,}4\%$ (60-24) & 3,6\% & $n$ de los más altos \\
EE.\,UU. & 18.992 & 66.857 & 2,0\% (60-24); 1,3\% (00-24) & 0,8\% & la frontera \\
\bottomrule
\end{tabular}
\end{center}

\begin{center}
\begin{tikzpicture}
\begin{axis}[
  width=\textwidth, height=6.5cm,
  ymode=log, log basis y=10,
  xmin=1960, xmax=2024, ymin=150, ymax=90000,
  xlabel={Año}, ylabel={$y$ (USD de 2015, log)},
  legend style={at={(0.5,-0.22)}, anchor=north, legend columns=5, font=\small, /tikz/every even column/.append style={column sep=6pt}},
  xticklabel style={/pgf/number format/1000 sep=},
  tick label style={font=\small}, label style={font=\small},
  unbounded coords=jump,
]
\addplot[thick, blue] table[x=anio,y=BWA]{datos-africa/ypc.dat};
\addlegendentry{Botsuana}
\addplot[thick, red] table[x=anio,y=COD]{datos-africa/ypc.dat};
\addlegendentry{RD del Congo}
\addplot[thick, orange, dashed] table[x=anio,y=RWA]{datos-africa/ypc.dat};
\addlegendentry{Ruanda}
\addplot[thick, gray, dotted] table[x=anio,y=NER]{datos-africa/ypc.dat};
\addlegendentry{Níger}
\addplot[thick, black] table[x=anio,y=USA]{datos-africa/ypc.dat};
\addlegendentry{EE.\,UU.}
\end{axis}
\end{tikzpicture}
\end{center}

\section*{Comentes}
\textit{Verdadero, Falso o Incierto. Justifique su respuesta.}

\begin{enumerate}
  \item En 1960, Botsuana y la RD del Congo eran mucho más pobres que EE.\,UU. Por lo tanto, según el modelo de Solow,
  ambos debieron crecer más rápido que EE.\,UU. desde entonces.
  \item Entre 1960 y 2024, la brecha entre EE.\,UU. y la RD del Congo pasó de 15 a 121 veces, y la de EE.\,UU. con Níger,
  de 26 a 114 veces. Esta evidencia refuta el modelo de Solow.
  \item Botsuana creció 9\% anual por persona entre 1966 y 1990, pero solo 1,3\% entre 2000 y 2024. Esa
  desaceleración muestra que su modelo de desarrollo fracasó.
  \item Desde el año 2000, Botsuana y EE.\,UU. crecen a la misma tasa por persona (1,3\%). Entonces, según Solow,
  ambos ya tienen el mismo producto por trabajador de estado estacionario.
  \item En 1994, tras el genocidio, el PIB de Ruanda cayó 50\%, y en 1995 creció 35\%. El modelo de Solow no puede
  explicar un crecimiento de 35\%, porque la tecnología no crece a esa tasa.
  \item Níger tiene uno de los crecimientos de población más altos del mundo ($n \approx 3{,}6\%$). Según el modelo
  de Solow con progreso tecnológico, eso explica que su PIB por persona crezca menos que el de EE.\,UU. en el largo
  plazo.
  \item Si la RD del Congo subiera su tasa de inversión del 12\% al 35\% (la de Botsuana durante su despegue),
  alcanzaría el producto por trabajador de estado estacionario de Botsuana. (Use $\alpha = 1/3$, $\delta = 0{,}05$
  y $g = 0{,}02$.)
\end{enumerate}

\vfill
\begin{center}\itshape --- La pauta empieza en la página siguiente ---\end{center}

\newpage
\section*{Pauta}

\textbf{1. Falso.} Eso es la \emph{convergencia absoluta}, que supone que todos tienen la misma función de producción
($A$, $h$) y los mismos $s$ y $n$, y por lo tanto el mismo estado estacionario. Solow predice convergencia
\emph{condicional}: cada país crece según su distancia a su propio $y_i^*$. Los datos lo muestran: Botsuana creció
9\% anual en 1966-90, y la RD del Congo, más rica que Botsuana en 1960 (3,3 veces), cayó 1,3\% anual. Ambos eran pobres, pero tenían
estados estacionarios muy distintos. El de Botsuana subió con $s \approx 35\%$ e instituciones que protegieron
la renta de los diamantes. El del Congo bajó con la guerra, una inversión baja y una $A$ en caída. Un país puede
ser pobre y crecer poco si ya está cerca de su propio $y^*$, que es bajo.

\textbf{2. Falso.} Refuta la convergencia \emph{absoluta}, no el modelo. Como en la clase 4 (el cociente
P90/P10 subió de $\sim$20 a $\sim$45), a nivel mundial la relación entre ser pobre y crecer rápido se rompe. Con
convergencia condicional, $\dfrac{y_{EEUU}^*}{y_{COD}^*} = \dfrac{A_{EEUU}h_{EEUU}}{A_{COD}h_{COD}}
\left(\dfrac{s_{EEUU}/(n_{EEUU}+g+\delta)}{s_{COD}/(n_{COD}+g+\delta)}\right)^{\frac{\alpha}{1-\alpha}}$ puede ser
enorme, y la brecha puede crecer si $A_{COD}$ cae (lo que ocurrió con el conflicto) o si $s_{COD}$ baja. Lo que Solow
\emph{sí} predice es que se crece más rápido cuanto más lejos se está del propio estado estacionario. Esto se ve
en clubes con $A$ y $h$ parecidos (la OCDE) y en casos como Corea o Botsuana.

\textbf{3. Falso.} Es exactamente lo que predice la dinámica de transición. Con $\dot{\kt}/\kt = s\,\kt^{\,\alpha-1}-(n+g+\delta)$,
cuando $\kt$ está lejos bajo $\kt^*$, la curva $s\kt^{\alpha-1}$ queda muy por encima de $n+g+\delta$, y $k$ e $y$ crecen
mucho más que $g$. A medida que $\kt \to \kt^*$, la brecha se cierra y $\dot{y}/y \to g$. Que Botsuana hoy crezca 1,3\%,
lo mismo que EE.\,UU., es la señal de que llegó cerca de su sendero de crecimiento balanceado: \emph{convergió}
a su $y^*$. No es un fracaso. \textit{Gráfico:} curva $s\kt^{\alpha-1}$ decreciente, recta $n+g+\delta$; la distancia
vertical es la tasa de crecimiento, y se achica al acercarse a $\kt^*$.

\textbf{4. Falso.} En el sendero de crecimiento balanceado, \emph{todos} crecen a $g$ (por persona), sin importar su nivel:
$\dot{y}/y = \dot{A}/A = g$. Crecer igual es un resultado de \emph{tasa}. El \emph{nivel} depende de
$A_i h_i\,(s_i/(n_i+g+\delta))^{\alpha/(1-\alpha)}$. Botsuana redujo la brecha de 47 veces (1966) a 9,6 veces (2024),
pero la brecha que queda es de nivel, por $A$, $h$, $s$ y $n$ distintos, y no se cierra sola. Es la diferencia entre
efecto nivel y efecto crecimiento.

\textbf{5. Falso.} Es dinámica de transición, no tecnología. El shock dejó a la economía muy por debajo de su estado
estacionario: capital destruido, producción interrumpida. Con $\dot{k}/k = s\,y/k - (n+\delta)$, si $k$ cae mucho,
$y/k$ sube (por rendimientos decrecientes, $PMg_K = \alpha y/k$ es alta) y la tasa de crecimiento se dispara por sobre $g$.
Ruanda recuperó su PIB por persona de 1993 en el año 2000 y creció 8,3\% anual en 1994-2004. Es el mismo argumento del
terremoto de Chile en 2010. \textit{Matiz que suma puntos:} el genocidio también destruyó población e instituciones
($L$ y $A$), así que parte del rebote es la recuperación de $A$, no solo de $k$.

\textbf{6. Falso (en Solow).} En Solow, $n$ tiene un \emph{efecto nivel}, no un efecto crecimiento: un $n$ mayor desplaza
la recta $n+g+\delta$ hacia arriba y \emph{baja} $\kt^*$ e $y^*$ (más trabajadores que equipar con el mismo ahorro). En el largo plazo,
$\dot{y}/y = g$ igual para Níger y para EE.\,UU. El bajo crecimiento observado de Níger ($-0{,}4\%$) se explica por un
$y^*$ bajo y por caídas de $A$, no por $n$ en sí. \textit{Para ir al 7:} en Romer semiendógeno es al revés,
$g_A = \lambda n/(1-\phi)$, y un $n$ mayor \emph{aumenta} el crecimiento de largo plazo (más gente que piensa, más ideas).
El profesor dijo que esto es «súper controversial», y mencionó India y países africanos: Níger es un contraejemplo.

\textbf{7. Falso (o Incierto).} En estado estacionario, con $\alpha/(1-\alpha) = 1/2$:
\[
\frac{y^*_{COD,\ s=35\%}}{y^*_{COD,\ s=12\%}} = \left(\frac{0{,}35}{0{,}12}\right)^{1/2} \approx 1{,}71 .
\]
Aun corrigiendo por $n$ (3,2\% en el Congo versus 1,7\% en Botsuana):
$\left(\dfrac{0{,}032+0{,}07}{0{,}017+0{,}07}\right)^{1/2} \approx 1{,}08$, $s$ y $n$ juntos explican un factor de $\approx 1{,}85$.
La brecha real es $6.952/554 \approx 12{,}5$. Queda un factor de $\approx 12{,}5/1{,}85 \approx 6{,}8$ que es $A\cdot h$: el
\emph{residuo de Solow}. Calza con la clase 4: la productividad explica entre la mitad y dos tercios de las diferencias de
ingreso. Subir $s$ tiene un efecto nivel acotado por los rendimientos decrecientes; sin subir $A$ y $h$, el Congo no
alcanza a Botsuana. (Incierto si se argumenta que una inversión mayor podría venir con mejores instituciones y subir $A$.)

\vfill
{\raggedright\small\textit{Fuente:} Banco Mundial, World Development Indicators (consultados el 6-11-2026): \texttt{NY.GDP.PCAP.KD},
\texttt{NY.GDP.MKTP.KD.ZG}, \texttt{SP.POP.GROW}, \texttt{NE.GDI.TOTL.ZS}. Las tasas son crecimientos compuestos calculados de esas series.\par}
\end{document}
